\chapter{Linux 内核数据结构与算法专题}

学习 Linux 内核数据结构，不能把它当成另一套刷题模板。内核里的链表、哈希链、红黑树、XArray、Maple Tree、RCU、伙伴系统和 LRU 链表，都是在同一个工程环境里长出来的：对象生命周期由内核控制，很多对象不能随便移动，内存分配可能发生在严格上下文中，读写并发非常频繁，某个对象还常常要同时挂进多个索引。理解这些约束以后，再回看力扣题，会更容易明白为什么一个高质量答案必须先问访问模式，而不是先背容器名字。

本章的阅读目标不是背某个内核版本的字段名。Linux 头文件和内部结构会随着版本变化，字段、辅助宏和封装层都可能调整；真正稳定的是设计原则：节点常嵌入业务对象，比较逻辑常由调用者提供，读路径尽量减少分配和锁竞争，删除对象常分成“从索引摘除”和“等读者离开后释放”两个阶段。读源码时请优先看接口语义、热路径和不变量，再看具体字段。

本章先讲内核为什么不直接使用 STL 容器，再依次分析侵入式链表、hlist、红黑树、XArray、Maple Tree、RCU、调度队列、内存管理和 I/O 队列。每一节都只围绕一个单点问题展开：这个结构解决什么访问模式，为什么普通刷题容器不够，关键不变量是什么，删除和并发为什么更难。最后的实验再把这些思想落回简化 \cpp 代码，帮助读者把源码阅读转成可复用的数据结构能力。

\begin{principlebox}{源码阅读路线}
\begin{tabular}{@{}p{0.24\linewidth}p{0.68\linewidth}@{}}
链表 & 读 \filepath{include/linux/list.h}，先看 \code{struct list_head}、初始化、插入、删除和遍历宏。\\
哈希链 & 读 \code{hlist} 相关接口，理解桶头和节点为什么分开，为什么节点里要保存前驱指针位置。\\
有序树 & 读 \filepath{include/linux/rbtree.h}，看节点嵌入、旋转由库完成、比较由调用者完成。\\
整数索引 & 读 \filepath{include/linux/xarray.h}，关注整数 ID 到指针的稀疏映射。\\
区间索引 & 读 \filepath{include/linux/maple_tree.h}，关注范围查询、VMA 这类区间对象为什么不适合普通数组。\\
并发读 & 读 \filepath{include/linux/rcupdate.h}，先理解读侧临界区、发布新版本和延迟释放。\\
\end{tabular}
\end{principlebox}

本章建议对照 Linux 官方文档阅读：链表见 \url{https://docs.kernel.org/core-api/list.html}，红黑树见 \url{https://docs.kernel.org/core-api/rbtree.html}，XArray 见 \url{https://docs.kernel.org/core-api/xarray.html}，Maple Tree 见 \url{https://docs.kernel.org/core-api/maple_tree.html}，RCU 入门见 \url{https://docs.kernel.org/RCU/whatisRCU.html}。书中代码是教学化简化版，用来解释不变量和访问模式，不代表内核源码的逐行复制。

\section{为什么内核不用 STL 容器}

在用户态 C++ 里，\code{std::vector}、\code{std::list}、\code{std::map} 和 \code{std::unordered_map} 会把元素所有权、分配器、迭代器、异常安全和泛型接口封装在容器内部。这个模型非常适合普通应用：你把元素放进容器，容器负责管理节点和内存，算法通过迭代器访问元素。内核的环境不一样。内核主要用 C 写成，不能依赖 C++ 异常和模板；很多路径不能随意阻塞分配；同一个对象常常同时属于多个结构；对象释放必须和锁、引用计数、RCU 宽限期配合。于是内核更偏向\term{侵入式数据结构}{intrusive data structure}：链接节点嵌进业务对象，数据结构只负责组织对象，不拥有对象。

以一个进程、一个页缓存对象或一个定时器对象为例，它可能同时需要按链表顺序被扫描，按哈希键被快速查找，按时间或地址范围被有序索引。如果每种索引都单独分配一个外部节点，再让节点指向业务对象，内存分配次数会增加，缓存局部性更差，释放路径也更复杂。侵入式结构让对象自己带着多个链接字段：一个字段挂进 LRU 链表，另一个字段挂进哈希桶，另一个字段挂进树。代价是调用者必须非常清楚节点属于哪个结构、何时插入、何时删除、对象何时还能被访问。

这和刷题中的 LRU 缓存有直接关系。普通写法是 \code{unordered_map<Key, list<Node>::iterator>} 加 \code{list<Node>}。哈希表负责按 key 定位，链表负责维护新旧顺序。内核场景常进一步把链表节点嵌到缓存对象内部，哈希索引也嵌到对象内部；对象本身不因为进入链表而被复制。你在面试中说“哈希表加双向链表”时，如果能补一句“哈希表负责定位，链表负责顺序，节点身份必须稳定”，就已经接近源码思维。

\section{侵入式链表：\code{list_head}}

Linux 的双向链表不是一个拥有元素的容器，而是一个很小的链接节点。典型形式可以抽象成下面这样：

\begin{lstlisting}[language=C++,caption={教学化简版：侵入式双向链表节点}]
struct ListHead {
    ListHead* next;
    ListHead* prev;
};

struct PageLike {
    int id;
    bool active;
    ListHead lru_node;
    ListHead dirty_node;
};
\end{lstlisting}

\code{PageLike} 里有两个链表节点，意味着同一个对象可以同时挂进两个不同链表。链表节点不保存 \code{id}，也不知道外层类型。遍历链表时，源码会通过成员地址反推外层对象地址，这就是常说的 \code{container_of} 思想：已知一个成员在结构体中的偏移，就能从成员指针减去偏移得到结构体起始地址。它不是魔法，而是 C 语言对象布局和偏移计算的组合。

\begin{principlebox}{图示：从节点回到对象}
\begin{tabular}{@{}p{0.2\linewidth}p{0.72\linewidth}@{}}
业务对象 & \codeesc{PageLike \{ id, active, lru\_node, dirty\_node \}}\\
链表视角 & 链表只看见一串 \code{ListHead}，不知道 \code{PageLike} 的其他字段。\\
反推对象 & 遍历到 \code{lru_node} 后，用成员偏移回到包含它的 \code{PageLike}。\\
设计收益 & 不额外分配链表节点，同一对象能进入多个链表，节点地址稳定。\\
\end{tabular}
\end{principlebox}

链表操作的核心不变量很简单：对任意节点 \code{x}，如果它在链表中，那么 \code{x->next->prev == x} 且 \code{x->prev->next == x}。插入节点就是同时改四条边；删除节点也是同时改四条边。错误通常不来自算法难，而来自顺序和生命周期：删除前是否还在链表中，删除后是否把节点重新初始化，遍历过程中删除当前节点是否提前保存了下一个节点。

\begin{lstlisting}[language=C++,caption={教学化简版：在两个节点之间插入新节点}]
void insert_between(ListHead* previous, ListHead* next, ListHead* node)
{
    node->next = next;
    node->prev = previous;
    previous->next = node;
    next->prev = node;
}

void remove_node(ListHead* node)
{
    ListHead* previous = node->prev;
    ListHead* next = node->next;
    previous->next = next;
    next->prev = previous;
    node->next = node;
    node->prev = node;
}
\end{lstlisting}

只给出两个函数还不够，读者必须看到“业务对象、嵌入节点、链表头、遍历和删除”如何组合起来。下面是一个完整教学版：\code{PageLike} 同时带有 \code{lru_node} 和 \code{dirty_node} 两个节点，\code{IntrusivePageList} 只操作其中一个成员。链表不复制 \code{PageLike}，也不释放 \code{PageLike}；它只维护对象之间的链接关系。

\begin{lstlisting}[language=C++,caption={侵入式双向链表的完整教学实现}]
struct ListHead {
    ListHead* next = nullptr;
    ListHead* prev = nullptr;
};

void init_list_head(ListHead* node)
{
    node->next = node;
    node->prev = node;
}

bool list_empty(const ListHead* head)
{
    return head->next == head;
}

void list_add_between(ListHead* previous, ListHead* next, ListHead* node)
{
    node->next = next;
    node->prev = previous;
    previous->next = node;
    next->prev = node;
}

void list_remove(ListHead* node)
{
    ListHead* previous = node->prev;
    ListHead* next = node->next;
    previous->next = next;
    next->prev = previous;
    init_list_head(node);
}

struct PageLike {
    int id = 0;
    bool active = false;
    ListHead lru_node;
    ListHead dirty_node;

    explicit PageLike(int page_id) : id(page_id)
    {
        init_list_head(&this->lru_node);
        init_list_head(&this->dirty_node);
    }
};

PageLike* page_from_lru_node(ListHead* node)
{
    auto* address = reinterpret_cast<std::byte*>(node) -
                    offsetof(PageLike, lru_node);
    return reinterpret_cast<PageLike*>(address);
}

class IntrusivePageList {
public:
    IntrusivePageList()
    {
        init_list_head(&this->head_);
    }

    bool empty() const
    {
        return list_empty(&this->head_);
    }

    void push_front(PageLike& page)
    {
        list_add_between(&this->head_, this->head_.next, &page.lru_node);
    }

    void push_back(PageLike& page)
    {
        list_add_between(this->head_.prev, &this->head_, &page.lru_node);
    }

    void move_to_front(PageLike& page)
    {
        list_remove(&page.lru_node);
        this->push_front(page);
    }

    PageLike* pop_back()
    {
        if (this->empty()) {
            return nullptr;
        }
        ListHead* node = this->head_.prev;
        list_remove(node);
        return page_from_lru_node(node);
    }

    std::vector<int> ids() const
    {
        std::vector<int> result;
        for (ListHead* node = this->head_.next;
             node != &this->head_;
             node = node->next) {
            result.push_back(page_from_lru_node(node)->id);
        }
        return result;
    }

private:
    ListHead head_;
};
\end{lstlisting}

这个实现里最重要的不是 \code{reinterpret_cast}，而是对象身份。链表节点只是 \code{PageLike} 的一个成员；节点地址稳定，外层对象就能稳定地被多个索引引用。若同一个 \code{PageLike} 要同时进入 LRU 链表和脏页链表，必须使用两个不同的 \code{ListHead} 成员，不能把同一个节点挂进两个链表。否则第二次插入会覆盖第一次插入的前后指针，两个链表都会被破坏。

这段代码要和力扣链表题联系起来看。反转链表强调“改 \code{next} 前保存后继”，K 个一组翻转强调“先确定分组边界”，LRU 强调“已知节点常数时间移动到头部”。内核链表把这些细节放大了：因为节点可能被很多模块引用，错误删除不是返回错误答案，而是破坏内核对象关系。刷题阶段写链表题，也应该养成明确前驱、当前、后继和已处理前缀的习惯。

\begin{labbox}
实验：写一个最小侵入式链表。定义 \codeesc{struct Item \{ int key; ListHead node; \}}，创建三个对象并把它们插入同一个链表。遍历时不要在 \code{ListHead} 里存 key，而是通过成员偏移回到 \code{Item}。再让 \code{Item} 增加第二个 \code{ListHead}，同时挂入另一个链表，观察同一对象如何拥有两种顺序。
\end{labbox}

\section{hlist：哈希桶里的瘦链表}

普通双向链表每个头节点也是一个 \code{list_head}，空链表时头节点的 \code{next} 和 \code{prev} 都指向自己。哈希表有大量桶，如果每个桶头都放完整双向节点，桶头本身会占用更多空间。Linux 的 \code{hlist} 可以理解为一种面向哈希桶优化的链表：桶头只保存第一个节点，节点保存下一个节点以及“指向当前节点指针的位置”。这样删除已知节点时仍然可以常数时间更新前驱链接，同时桶头更轻。

这背后的算法思想和刷题哈希表相同：哈希函数把 key 映射到桶，桶内用链表处理冲突。区别在于内核更关心节点分配、缓存局部性和并发删除。刷题里你使用 \code{unordered_map} 时，桶和节点都由标准库隐藏；读内核源码时，桶数组、节点链接和删除顺序都暴露在接口语义中。

\begin{principlebox}{哈希链的不变量}
\begin{tabular}{@{}p{0.22\linewidth}p{0.7\linewidth}@{}}
桶定位 & 哈希值决定先进入哪个桶，桶只缩小候选范围，不保证桶内有序。\\
冲突处理 & 同桶节点继续按 key 比较，不能只比较哈希值。\\
删除节点 & 已知节点删除要能改前驱的 \code{next}，否则必须从桶头重新扫描。\\
并发风险 & 读者遍历时删除节点，需要锁、引用计数或 RCU 保护生命周期。\\
\end{tabular}
\end{principlebox}

下面的教学版 hlist 用 C++ 写出核心结构。桶头 \code{HListHead} 只有一个 \code{first} 指针；节点 \code{HListNode} 保存 \code{next}，还保存 \code{previous_link}，后者指向“当前节点是从哪里被指到的”。如果节点是桶中第一个元素，\code{previous_link} 指向桶头的 \code{first}；如果节点前面还有节点，\code{previous_link} 指向前一个节点的 \code{next}。删除时只要改 \code{*previous_link}，就能把前驱链接直接接到后继。

\begin{lstlisting}[language=C++,caption={hlist 风格哈希桶：完整插入、查找和已知节点删除}]
struct HListNode {
    HListNode* next = nullptr;
    HListNode** previous_link = nullptr;
};

struct HListHead {
    HListNode* first = nullptr;
};

bool hlist_unhashed(const HListNode* node)
{
    return node->previous_link == nullptr;
}

void hlist_add_head(HListHead* head, HListNode* node)
{
    HListNode* first = head->first;
    node->next = first;
    if (first != nullptr) {
        first->previous_link = &node->next;
    }
    head->first = node;
    node->previous_link = &head->first;
}

void hlist_del(HListNode* node)
{
    if (hlist_unhashed(node)) {
        return;
    }
    HListNode* next = node->next;
    *node->previous_link = next;
    if (next != nullptr) {
        next->previous_link = node->previous_link;
    }
    node->next = nullptr;
    node->previous_link = nullptr;
}

struct KernelObject {
    int key = 0;
    std::string value;
    HListNode hash_node;
};

KernelObject* object_from_hash_node(HListNode* node)
{
    auto* address = reinterpret_cast<std::byte*>(node) -
                    offsetof(KernelObject, hash_node);
    return reinterpret_cast<KernelObject*>(address);
}

class IntrusiveHashTable {
public:
    explicit IntrusiveHashTable(int bucket_count)
        : buckets_(static_cast<std::size_t>(bucket_count))
    {
        if (bucket_count <= 0) {
            throw std::invalid_argument("bucket_count");
        }
    }

    void insert(KernelObject& object)
    {
        if (!hlist_unhashed(&object.hash_node)) {
            hlist_del(&object.hash_node);
        }
        HListHead& bucket = this->bucket_for(object.key);
        hlist_add_head(&bucket, &object.hash_node);
    }

    KernelObject* find(int key) const
    {
        const HListHead& bucket = this->bucket_for(key);
        HListNode* node = bucket.first;
        while (node != nullptr) {
            KernelObject* object = object_from_hash_node(node);
            if (object->key == key) {
                return object;
            }
            node = node->next;
        }
        return nullptr;
    }

    bool erase(int key)
    {
        KernelObject* object = this->find(key);
        if (object == nullptr) {
            return false;
        }
        hlist_del(&object->hash_node);
        return true;
    }

private:
    HListHead& bucket_for(int key)
    {
        const auto index =
            std::hash<int>{}(key) % static_cast<std::size_t>(this->buckets_.size());
        return this->buckets_[index];
    }

    const HListHead& bucket_for(int key) const
    {
        const auto index =
            std::hash<int>{}(key) % static_cast<std::size_t>(this->buckets_.size());
        return this->buckets_[index];
    }

    std::vector<HListHead> buckets_;
};
\end{lstlisting}

这个实现故意不让 \code{IntrusiveHashTable} 拥有 \code{KernelObject}，因为内核风格的数据结构通常只是索引对象。删除 \code{erase} 只把对象从哈希桶摘掉，不负责释放对象；释放必须由对象生命周期管理者决定。若有并发读者正在遍历桶，\code{hlist_del} 之后也不能立刻 \code{delete} 对象，这正是后面 RCU 要解决的问题。刷题里的 \code{unordered_map} 把这些细节藏起来；源码学习则必须把“摘除索引”和“释放对象”分清。

把它映射到面试题，两数之和和前缀和计数都在使用“key 到历史状态”的映射。高质量题解不应只写 \code{unordered_map}，还要说明 key 是什么、value 是什么、什么时候查询、什么时候插入、是否允许当前元素和自己匹配。内核里的 hlist 把这些语义拆得更底层：桶只是入口，节点才是业务对象，生命周期由调用者控制。

\section{红黑树：有序索引和调用者比较}

红黑树解决的是动态有序集合问题：插入、删除、查找、前驱、后继都能在对数复杂度内完成。C++ 的 \code{std::map} 把 key、value、比较器和树节点封装起来；Linux 的 \code{rbtree} 更底层，树节点通常嵌入业务对象，库函数负责旋转和颜色维护，但“往左还是往右”的比较逻辑由调用者写。这样做让内核对象可以按自己的字段组织成树，也避免为每个节点额外包装一个通用容器节点。

\begin{lstlisting}[language=C++,caption={教学化简版：业务对象嵌入树节点}]
struct RbNode {
    RbNode* left;
    RbNode* right;
    RbNode* parent;
    bool red;
};

struct TimerLike {
    std::int64_t expires;
    RbNode tree_node;
};
\end{lstlisting}

红黑树的细节可以很复杂，但教材阶段先抓三件事。第一，它是二叉搜索树，左子树 key 小，右子树 key 大。第二，它通过颜色约束和旋转保持近似平衡，避免退化成链。第三，插入和删除之后，局部旋转不会破坏中序顺序，只是在保持排序的前提下修复高度和颜色约束。刷题中的“搜索旋转排序数组”和这里的旋转不是同一件事；树旋转是局部结构调整，中序序列不变。

\begin{principlebox}{红黑树和刷题有序容器}
\begin{tabular}{@{}p{0.24\linewidth}p{0.68\linewidth}@{}}
需要顺序 & 只要题目问前驱、后继、最小大于、最大小于，哈希表就不够。\\
动态更新 & 如果元素会插入删除，排序数组每次移动成本高，有序树更合适。\\
比较责任 & 内核调用者自己决定 key；C++ \code{map} 把比较器藏进容器类型。\\
正确性 & 树旋转保持中序顺序，所以查找语义不变，只改善平衡。\\
\end{tabular}
\end{principlebox}

在内核中，有序树常用于按时间、地址或其他可比较 key 组织对象。具体模块会随版本演进而变化，例如某些历史结构可能从红黑树迁移到 Maple Tree 或其他结构。读源码时不要死记“某模块一定使用某树”，而要问：这个模块是否需要动态有序查找，是否需要范围查询，是否需要前驱后继，是否需要频繁插入删除。

\section{XArray：稀疏整数索引}

很多内核对象需要用整数 ID 查找：页缓存按页索引查找，文件描述符按编号查找，设备和 ID 分配也常见“整数到对象”的映射。直接开一个巨大数组很浪费，因为 ID 空间可能稀疏；用普通哈希表能查找，但范围扫描、按序遍历和某些并发语义不一定理想。XArray 可以理解为一种面向无符号长整数索引的稀疏数组结构，它以树状分层方式把大整数索引拆开，按需分配中间节点。

从算法上看，这和 Trie 有相似味道：字符串 Trie 每层消耗一个字符，XArray 这类整数索引结构每层消耗索引的一部分位。若索引空间很大但实际对象不多，树只为存在对象的路径分配节点。若需要查找某个 ID，沿着索引位段向下走；若需要扫描一段 ID，结构可以按顺序推进。它比“巨大数组”省空间，比“无序哈希”更重视索引顺序和内核并发接口。

\begin{principlebox}{从力扣题理解 XArray}
\begin{tabular}{@{}p{0.24\linewidth}p{0.68\linewidth}@{}}
数组 & 下标连续且范围小，直接数组最快。\\
哈希表 & key 稀疏且只需要等值查找，哈希表简单。\\
Trie / 分层索引 & key 可以按位或按字符分层，适合共享前缀和稀疏空间。\\
XArray & 面向内核对象的整数到指针映射，还要考虑遍历、标记和并发读写。\\
\end{tabular}
\end{principlebox}

刷题中类似的思想出现在前缀树、位 Trie、稀疏表和坐标压缩。区别是刷题往往只考虑单线程和答案复杂度，内核还要考虑节点分配、锁、RCU、标记位和对象生命周期。如果你读 XArray 感到细节很多，先不要陷进每个宏，先回答三个问题：索引如何分层，空洞如何表示，对象指针何时对读者可见。

\section{Maple Tree：范围映射和 VMA 思维}

普通映射结构处理的是点查询：给定 key 找对象。内存管理中大量问题是范围查询：给定地址，找覆盖它的虚拟内存区域；插入一个新区间，需要检查与相邻区间是否重叠；删除或分裂区间后，还要保持后续查询高效。Maple Tree 是 Linux 中用于范围映射的一类重要结构，现代内核用它承载诸如虚拟内存区域管理这类范围索引场景。你不需要在刷题里复刻它，但要理解为什么“区间对象”比“单点 key”更复杂。

力扣里的合并区间、插入区间、删除被覆盖区间和日程安排，都是范围结构的简化版。排序数组适合离线一次性处理；如果区间会动态插入、删除、查询，就需要有序结构维护相邻关系。Maple Tree 的动机也是如此：地址空间很大，区间数量动态变化，查询频繁，且读路径性能很关键。它的实现细节比普通红黑树复杂，但读者先抓住范围映射语义即可。

\begin{principlebox}{区间结构的单点分析}
\begin{tabular}{@{}p{0.24\linewidth}p{0.68\linewidth}@{}}
点查询 & 给定地址，找到起点不大于它且终点覆盖它的区间。\\
插入 & 新区间必须和前后相邻区间比较，判断重叠、合并或拒绝。\\
删除 & 可能删除整个区间，也可能把一个区间切成左右两段。\\
遍历 & 常需要按地址顺序扫描下一段区域，不能只支持等值查找。\\
\end{tabular}
\end{principlebox}

学习这部分时，建议先用 \code{std::map<int, Interval>} 写一个小型区间表：key 是区间起点，value 是终点和属性。实现 \code{find(point)}、\code{insert(l, r)}、\code{erase(l, r)} 三个操作。写完以后再问自己：如果区间数量巨大、读多写少、对象释放有并发读者，这个简单 map 还缺什么？这个问题会自然引出内核结构的复杂性。

\begin{lstlisting}[language=C++,caption={动态区间表：点查询、插入和删除子区间}]
struct MemoryArea {
    std::int64_t end = 0;
    std::string permission;
};

class VirtualMemoryMap {
public:
    const MemoryArea* find(std::int64_t address) const
    {
        const auto after = this->areas_.upper_bound(address);
        if (after == this->areas_.begin()) {
            return nullptr;
        }
        auto current = after;
        --current;
        if (current->second.end <= address) {
            return nullptr;
        }
        return &current->second;
    }

    bool insert(std::int64_t begin,
                std::int64_t end,
                std::string permission)
    {
        if (begin >= end) {
            return false;
        }
        const auto next = this->areas_.lower_bound(begin);
        if (next != this->areas_.end() && next->first < end) {
            return false;
        }
        if (next != this->areas_.begin()) {
            auto previous = next;
            --previous;
            if (previous->second.end > begin) {
                return false;
            }
        }
        this->areas_.emplace(begin, MemoryArea{end, std::move(permission)});
        return true;
    }

    bool erase(std::int64_t begin, std::int64_t end)
    {
        if (begin >= end) {
            return false;
        }
        auto current = this->areas_.upper_bound(begin);
        if (current == this->areas_.begin()) {
            return false;
        }
        --current;
        if (current->second.end < end) {
            return false;
        }

        const std::int64_t old_begin = current->first;
        const MemoryArea old_area = current->second;
        this->areas_.erase(current);

        if (old_begin < begin) {
            this->areas_.emplace(old_begin,
                                 MemoryArea{begin, old_area.permission});
        }
        if (end < old_area.end) {
            this->areas_.emplace(end,
                                 MemoryArea{old_area.end, old_area.permission});
        }
        return true;
    }

    std::vector<std::pair<std::int64_t, MemoryArea>> snapshot() const
    {
        std::vector<std::pair<std::int64_t, MemoryArea>> result;
        result.reserve(this->areas_.size());
        for (const auto& entry : this->areas_) {
            result.push_back(entry);
        }
        return result;
    }

private:
    std::map<std::int64_t, MemoryArea> areas_;
};
\end{lstlisting}

这段代码把区间不变量写得很明确：所有区间按起点排序，且相邻区间互不重叠。点查询时，真正可能覆盖 \code{address} 的只会是起点不大于它的最后一个区间；插入时，只需要检查后继是否从新区间内部开始，再检查前驱是否延伸到新区间内部；删除时，必须支持从原区间左侧、右侧或中间切掉一段。若从中间删除，原区间会被分裂成左右两个区间。这正是内存 VMA 类问题比普通 key-value 查询更难的地方。

\section{RCU：读多写少下的数据结构生命周期}

RCU 的全称是 Read-Copy-Update。它不是某个容器，而是一种并发读写思想：读者在读侧临界区内尽量不加重锁；写者创建或修改新版本，把新指针发布给后来的读者；旧对象不能立刻释放，必须等所有可能还看见旧版本的读者离开之后再回收。这个“延迟释放”会直接改变数据结构删除操作的含义。

普通单线程链表删除可以理解为两步：改前后指针、释放节点。RCU 场景下必须拆得更细：第一步从索引中摘除，让新读者找不到它；第二步等待宽限期，让旧读者不再持有它；第三步才释放内存。如果把摘除和释放混成一步，读者可能拿着旧指针访问已经释放的对象。很多内核数据结构的复杂性，都来自这类生命周期问题。

\begin{principlebox}{RCU 删除的三阶段}
\begin{tabular}{@{}p{0.22\linewidth}p{0.7\linewidth}@{}}
逻辑删除 & 从链表、树或哈希表中摘除，使后续查找不再得到该对象。\\
等待读者 & 等待读侧临界区结束，确保旧读者不再访问旧对象。\\
物理释放 & 释放内存或归还缓存，此时对象生命周期才真正结束。\\
\end{tabular}
\end{principlebox}

这和刷题中的“懒删除”有相似但不相同的味道。滑动窗口中位数用堆时，过期元素先标记，等它到堆顶再弹出；RCU 中对象先从新索引中不可达，再等旧读者离开后释放。二者都说明删除不一定等于立刻销毁，区别在于刷题懒删除主要为了弥补容器能力，RCU 延迟释放主要为了并发安全。

\begin{lstlisting}[language=C++,caption={RCU 思想教学版：复制、发布和延迟释放}]
struct RoutingTable {
    std::unordered_map<int, std::string> route_by_id;
};

class RcuReadGuard {
public:
    explicit RcuReadGuard(std::atomic<int>& readers) : readers_(readers)
    {
        this->readers_.fetch_add(1, std::memory_order_acquire);
    }

    RcuReadGuard(const RcuReadGuard&) = delete;
    RcuReadGuard& operator=(const RcuReadGuard&) = delete;

    ~RcuReadGuard()
    {
        this->readers_.fetch_sub(1, std::memory_order_release);
    }

private:
    std::atomic<int>& readers_;
};

class RcuRoutingTable {
public:
    RcuRoutingTable()
        : current_(std::make_shared<const RoutingTable>())
    {
    }

    std::optional<std::string> lookup(int id) const
    {
        RcuReadGuard guard(this->reader_count_);
        std::shared_ptr<const RoutingTable> table =
            std::atomic_load_explicit(&this->current_, std::memory_order_acquire);
        const auto found = table->route_by_id.find(id);
        if (found == table->route_by_id.end()) {
            return std::nullopt;
        }
        return found->second;
    }

    void update(int id, std::string route)
    {
        std::shared_ptr<const RoutingTable> old_table =
            std::atomic_load_explicit(&this->current_, std::memory_order_acquire);
        auto next_table = std::make_shared<RoutingTable>(*old_table);
        next_table->route_by_id[id] = std::move(route);

        std::atomic_store_explicit(
            &this->current_,
            std::static_pointer_cast<const RoutingTable>(next_table),
            std::memory_order_release);
        this->retired_.push_back(std::move(old_table));
        this->reclaim_when_quiescent();
    }

    void erase(int id)
    {
        std::shared_ptr<const RoutingTable> old_table =
            std::atomic_load_explicit(&this->current_, std::memory_order_acquire);
        auto next_table = std::make_shared<RoutingTable>(*old_table);
        next_table->route_by_id.erase(id);

        std::atomic_store_explicit(
            &this->current_,
            std::static_pointer_cast<const RoutingTable>(next_table),
            std::memory_order_release);
        this->retired_.push_back(std::move(old_table));
        this->reclaim_when_quiescent();
    }

    int retired_count() const
    {
        return static_cast<int>(this->retired_.size());
    }

private:
    void reclaim_when_quiescent()
    {
        if (this->reader_count_.load(std::memory_order_acquire) != 0) {
            return;
        }
        this->retired_.clear();
    }

    mutable std::atomic<int> reader_count_{0};
    std::shared_ptr<const RoutingTable> current_;
    std::vector<std::shared_ptr<const RoutingTable>> retired_;
};
\end{lstlisting}

这不是内核 RCU 的实现，只是把语义拆开给读者看。读者进入读侧临界区后读取当前快照；写者不原地修改旧表，而是复制出新表、修改副本、再发布新指针。旧表进入 \code{retired_}，只有当读者计数归零时才清理。真实内核不会用 \code{shared_ptr} 管理这些对象，也不会用一个全局计数粗糙模拟宽限期；但教学目标是相同的：删除不是立刻释放，生命周期必须覆盖所有可能仍在读取旧版本的读者。

\section{调度器：从队列到虚拟运行时间}

操作系统调度要解决的问题是：大量可运行任务中，下一个应该让谁运行，运行多久，如何兼顾公平性、响应性和吞吐。最简单的模型是队列：先来先服务或轮转调度。但真实系统里任务权重不同、交互需求不同、睡眠和唤醒频繁，单纯 FIFO 不够。Linux 的完全公平调度思想可以粗略理解为维护任务的虚拟运行时间，优先选择虚拟运行时间较小的任务，让每个任务按权重获得相对公平的 CPU 时间。

这个模型和算法题里的“动态最小值”有关。你可以把可运行任务看成候选集合，每次取虚拟运行时间最小者运行，运行后它的虚拟运行时间增加，再放回候选集合。若只需要取最小值，堆似乎可行；若还需要删除任意任务、按顺序遍历、处理权重和复杂的队列状态，有序树或更专门的数据结构会更适合。具体内核实现会随调度类变化，但核心问题始终是动态候选集合和公平性指标。

\begin{labbox}
实验：写一个简化调度模拟器。每个任务有 \code{id}、\code{weight}、\code{vruntime} 和剩余运行时间。用 \code{std::set} 按 \code{vruntime} 排序，每次取最小者运行一个时间片，再更新 \code{vruntime} 放回。然后把 \code{std::set} 换成 \code{priority_queue}，尝试支持“任务睡眠时从队列删除”，比较两种结构的接口差异。
\end{labbox}

\section{伙伴系统、SLUB 和 LRU：内存管理里的算法}

内存管理是数据结构密度很高的地方。伙伴系统把物理页按 2 的幂次阶组织成空闲块：需要分配小块时，从更大块拆分；释放时，如果相邻伙伴也空闲，就合并成更大块。这个算法的目标不是最优装箱，而是在可接受碎片下快速分配和释放连续页。它和刷题里的区间合并、二进制分解、堆式层级都有联系，但更强调页框编号、阶数和合并条件。

\begin{principlebox}{伙伴系统的思维模型}
\begin{tabular}{@{}p{0.22\linewidth}p{0.7\linewidth}@{}}
阶数 & 第 \code{k} 阶块包含 \code{2^k} 个连续页。\\
拆分 & 没有合适小块时，从更高阶块拆成两个伙伴。\\
合并 & 释放后若伙伴同阶且空闲，就合并成高一阶块。\\
不变量 & 每个空闲块只出现在它当前阶数对应的空闲链表中。\\
\end{tabular}
\end{principlebox}

小对象分配又是另一层问题。频繁为内核对象调用通用页分配器成本高，也会造成碎片。SLAB/SLUB 这类分配器把同类型对象放进缓存，从页中切出固定大小对象，复用已经初始化过的对象槽位。刷题里你通常不关心对象分配成本，但工程里大量小节点会让哈希表、链表和树的常数差异变得很明显。理解 SLUB 后，再看 \code{list} 节点分配、\code{unordered_map} 节点分配，就会知道“理论 O(1)”背后还有分配器和缓存局部性。

页面回收中的 LRU 思想也值得和力扣 LRU 缓存比较。力扣 LRU 通常要求 \code{get} 和 \code{put} 都是 O(1)，结构是哈希表加双向链表。内核页面回收更复杂：页面可能分为活跃和非活跃，文件页和匿名页也可能有不同链表，访问位、脏页、写回和引用计数都会影响回收判断。它不是一个简单的“最久未使用就删”，而是在可测量近似、并发成本和 I/O 成本之间折中。教材阶段要学到的不是所有策略细节，而是多链表如何表达状态分层。

\begin{lstlisting}[language=C++,caption={简化伙伴分配器：按阶拆分和合并}]
class BuddyAllocator {
public:
    explicit BuddyAllocator(int max_order) : free_by_order_(max_order + 1)
    {
        if (max_order < 0 || max_order > MAX_SUPPORTED_ORDER) {
            throw std::invalid_argument("max_order");
        }
        this->max_order_ = max_order;
        this->free_by_order_[static_cast<std::size_t>(max_order)].insert(0);
    }

    std::optional<int> allocate(int order)
    {
        if (order < 0 || order > this->max_order_) {
            return std::nullopt;
        }

        int current_order = order;
        while (current_order <= this->max_order_ &&
               this->free_by_order_[static_cast<std::size_t>(current_order)].empty()) {
            ++current_order;
        }
        if (current_order > this->max_order_) {
            return std::nullopt;
        }

        int start = *this->free_by_order_[static_cast<std::size_t>(current_order)].begin();
        this->free_by_order_[static_cast<std::size_t>(current_order)].erase(start);

        while (current_order > order) {
            --current_order;
            const int buddy = start + block_size(current_order);
            this->free_by_order_[static_cast<std::size_t>(current_order)].insert(buddy);
        }
        this->allocated_[start] = order;
        return start;
    }

    bool release(int start)
    {
        const auto found = this->allocated_.find(start);
        if (found == this->allocated_.end()) {
            return false;
        }

        int order = found->second;
        this->allocated_.erase(found);
        int current_start = start;

        while (order < this->max_order_) {
            const int buddy = current_start ^ block_size(order);
            auto& free_list = this->free_by_order_[static_cast<std::size_t>(order)];
            const auto buddy_it = free_list.find(buddy);
            if (buddy_it == free_list.end()) {
                break;
            }
            free_list.erase(buddy_it);
            current_start = std::min(current_start, buddy);
            ++order;
        }

        this->free_by_order_[static_cast<std::size_t>(order)].insert(current_start);
        return true;
    }

    std::vector<int> free_counts() const
    {
        std::vector<int> counts;
        counts.reserve(this->free_by_order_.size());
        for (const auto& free_list : this->free_by_order_) {
            counts.push_back(static_cast<int>(free_list.size()));
        }
        return counts;
    }

private:
    static int block_size(int order)
    {
        return 1 << order;
    }

    static constexpr int MAX_SUPPORTED_ORDER = 30;
    int max_order_ = 0;
    std::vector<std::set<int>> free_by_order_;
    std::unordered_map<int, int> allocated_;
};
\end{lstlisting}

这个实现用 \code{std::set<int>} 保存每一阶空闲块的起始页号，是为了让教学代码更容易检查伙伴是否空闲。真实内核会用更贴近页结构和链表的表示。分配时，如果目标阶没有空闲块，就向更高阶借一个块，然后逐层拆成两个伙伴，把右伙伴放回低一阶空闲表；释放时，用 \code{start xor 2^order} 算出同阶伙伴起点，若伙伴也空闲，就删除伙伴并合并到高一阶。这个模型能讲清两个核心事实：总空闲页数足够不代表有足够大的连续块，释放时能否合并取决于伙伴是否同阶且空闲。

\section{网络、I/O 和定时器中的队列结构}

系统里还有大量队列结构。网络包需要排队、分类、限速和调度；块 I/O 请求需要合并相邻请求、排序或分发到设备队列；定时器需要按到期时间组织事件；工作队列需要把任务从中断或提交路径转移到可睡眠上下文执行。每种队列都不是“先进先出”四个字能解释完的，关键仍然是访问模式：是否只从头出队，是否需要按优先级出队，是否需要按时间查找最早事件，是否需要取消任意节点。

这部分和刷题题型也能对应。普通 BFS 队列只需要 FIFO；0-1 BFS 需要双端队列，因为零权边放队首、一权边放队尾；Dijkstra 需要按当前距离取最小，适合堆；任务调度器可能需要最大频次、冷却时间和时间推进；区间调度需要按结束时间排序。真实内核把这些问题放在更严格的并发和生命周期约束下，但算法选择的出发点完全一致。

\section{如何把内核源码读成能力}

读内核源码最怕两种极端：一种是只看博客结论，觉得“内核用了红黑树”就结束；另一种是一打开头文件就追每个宏，最后忘了自己为什么读。正确方法是单点问题单点读。比如今天只读链表删除，就只回答：删除已知节点要改几条边，遍历时删除为什么要 safe 版本，删除后节点是否重新初始化。明天只读红黑树插入，就只回答：调用者如何找到插入位置，库函数如何链接节点，修复函数为什么不改变中序顺序。

\section{五个源码实验案例}

这一节把 Linux 数据结构压成五个可以动手的实验。实验不要求完全复刻内核，只要求把源码里的访问模式和不变量跑出来。每个实验都要写报告：对象是什么，索引是什么，插入、查询、删除分别维护哪个不变量，和对应力扣题有什么相同点和不同点。

\subsection{实验一：侵入式 LRU}

从力扣 146 LRU 缓存开始。普通答案通常是 \code{list} 加 \code{unordered_map}，链表节点由 \code{list} 管理，哈希表保存迭代器。侵入式版本改成对象自己带两个字段：key、value 和双向链表节点。哈希表从 key 映射到对象指针，链表只负责顺序。这样能直观看到内核设计的第一原则：对象不因为进入链表而被复制，链表也不拥有对象。

实验步骤如下。第一，写一个 \code{CacheEntry}，内部包含 \code{ListHead lru}。第二，写 \code{move_to_front}、\code{remove_tail}、\code{insert_front} 三个链表操作。第三，用 \code{unordered_map<int, CacheEntry*>} 定位对象。第四，容量满时先从链表尾部找到对象，再从哈希表删除 key，最后释放对象。报告里要特别说明删除顺序：如果先释放对象，再从哈希表或链表删除，就会留下悬空指针。

这个实验对应内核里的 LRU 思维，但不能简单说“内核也是 LRU 缓存”。真实页面回收会考虑活跃、非活跃、脏页、写回、引用和 NUMA 等因素；实验只保留最核心的索引组合：哈希表负责定位，双向链表负责顺序，节点身份必须稳定。

\subsection{实验二：哈希桶和 hlist 删除}

写一个小哈希表，桶数组大小固定为 16，每个桶保存一条链。第一版用普通单链表，只在节点里保存 \code{next}。当你拿到一个节点指针并想删除它时，必须从桶头重新扫描找到前驱。第二版让节点额外保存“指向当前节点的前驱链接位置”，也就是可以修改前驱的 \code{next} 或桶头的 \code{first}。这样已知节点删除就不需要重新扫描。

报告要回答三个问题。第一，为什么桶头不一定需要完整双向节点？第二，为什么节点保存前驱链接位置可以支持常数时间删除？第三，如果另一个读者正在遍历这个桶，删除节点后什么时候才能释放对象？第三个问题会自然连接到 RCU：摘除只是让新查找找不到它，释放还要等旧读者离开。

这个实验对应的力扣题不是某一道固定题，而是一类哈希题：两数之和、前缀和计数、字母异位词分组、四数相加 II。刷题里桶结构被标准库隐藏；实验把桶、冲突和删除成本显式暴露出来。

\subsection{实验三：区间表和 Maple Tree 思维}

用 \code{std::map<std::int64_t, Interval>} 写一个简化虚拟地址区间表。key 是区间起点，value 保存终点和权限字符串。实现三个操作：\code{find(addr)} 找覆盖某个地址的区间；\code{insert(left, right)} 插入不重叠区间；\code{erase(left, right)} 删除一个子区间，必要时把原区间切成左右两段。

这个实验的关键不是 map 代码，而是区间不变量。任意两个相邻区间必须不重叠，查询地址时只需要找起点不大于它的最后一个区间，再判断终点是否覆盖。插入时只需要检查前驱和后继；删除时要处理四种情况：删除整段、删除左侧、删除右侧、从中间挖掉一段。每种情况都要写一个用例。

Maple Tree 比这个实验复杂得多，但问题同源：地址空间是稀疏范围集合，查询频繁，插入删除会改变相邻关系。力扣 56、57、435、986 是离线区间处理；这个实验是动态区间索引。两者的桥梁是“排序后相邻区间决定合法性”。

\subsection{实验四：简化伙伴分配器}

设总页数为 \code{2^{10}}，每个空闲块用起始页号和阶数表示。维护 11 个空闲链表，第 \code{k} 个链表保存大小为 \code{2^k} 的空闲块。分配 \code{k} 阶块时，如果当前阶没有空闲块，就向上找更大阶，逐层拆分；释放时，计算伙伴块编号，如果伙伴同阶且空闲，就从空闲链表删除伙伴并合并成高一阶。

这个实验训练的是二进制边界。伙伴块的起始页号不是随便算的，它和阶数对齐；两个伙伴合并后起点取较小者，阶数加一。报告要打印每次操作后各阶空闲块数量，并给出一个碎片案例：虽然总空闲页数足够，但没有足够大的连续块，导致高阶分配失败。

对应的算法题包括区间合并、位运算对齐、堆式层级和内存池模拟。它们都说明一个原则：分配器不是只问总量，还问连续性、对齐和合并条件。

\subsection{实验五：RCU 风格读写模拟}

写一个读多写少的配置表。读者只读一个全局指针指向的不可变表；写者复制旧表，修改副本，然后原子发布新指针。旧表不要立即释放，而是放入延迟释放列表。为了模拟宽限期，可以让每个读者进入读侧临界区时增加计数，离开时减少计数；当读者计数为零时，写者再释放延迟列表里的旧表。

这个实验会让你看到 RCU 和普通锁的不同。普通互斥锁让读者和写者互相等待，逻辑简单但读路径重；RCU 让读者轻量，写者承担复制、发布和延迟释放成本。数据结构上最大的变化是删除不等于释放：先让新读者不可达，再等旧读者结束，最后释放对象。

把它和刷题懒删除比较。滑动窗口中位数里，堆中旧元素暂时不删，是因为 \code{priority_queue} 不支持删除任意元素；RCU 旧对象暂时不释放，是因为并发读者可能仍然持有旧指针。相同点是“逻辑不可见”和“物理释放”分离，不同点是目的和安全边界完全不同。

\begin{exercisebox}
本章作业：

1. 用 C++ 写一个侵入式链表小实验，要求同一个对象同时挂入两个链表，并说明两个链表节点字段分别代表什么。

2. 用 \code{std::map} 写一个区间表，支持点查询、插入不重叠区间、删除一个子区间。完成后解释它和 Maple Tree 解决的问题有什么相同点，工程约束有什么不同。

3. 写一个简化伙伴系统，只管理 \code{2^10} 个页框，支持按阶分配和释放。每次操作后打印各阶空闲链表长度，观察拆分和合并。

4. 写一个 RCU 思想模拟：读者线程只读取共享指针，写者复制新对象并发布，旧对象放进延迟释放列表。即使不真正实现内核 RCU，也要把“摘除”和“释放”拆成两步。

5. 选一道力扣题和一个内核结构做对照。例如 LRU 缓存对照页面回收链表，合并区间对照 VMA 区间管理，前 K 高频对照动态候选极值。要求写出相同点和不同点，不能只写“都是链表”。
\end{exercisebox}

\begin{principlebox}{本章复盘}
Linux 内核数据结构不是为了展示技巧，而是为访问模式、生命周期、并发和内存分配服务。侵入式链表强调对象身份稳定，hlist 强调哈希桶空间和删除效率，红黑树强调动态有序索引，XArray 强调稀疏整数映射，Maple Tree 强调范围查询，RCU 强调读多写少下的延迟释放。刷题时把这些思想降维使用，就是高质量题解里的状态职责、不变量和容器选择。
\end{principlebox}
